\chapter{Einleitung}\label{ch:einleitung}


\section{Problemstellung und Motivation}

\acp{llm} finden in der Praxis und insbesondere in Unternehmen zunehmend Anwendung. Diese Entwicklung lässt sich in zahlreichen Statistiken der vergangenen 
Jahre verfolgen. So stellt das Statistische Bundesamt fest, dass der Anteil der Unternehmen, die \ac{ki} einsetzen, im Zeitraum von 2021 bis 2025 von 
11\,\% auf 26\,\% gestiegen ist. Dabei ist zu beachten, dass Unternehmen aller Größenklassen in diese Statistik einbezogen wurden, auch solche mit 10 bis 
49 Beschäftigten. Bei Unternehmen mit 250 und mehr Beschäftigten stieg der Anteil im selben Zeitraum von 31\,\% auf 57\,\%~\cite{e01}. \\
Das ifo Institut verzeichnet für Mai 2026, dass 54,5\,\% der Unternehmen in Deutschland \ac{ki} verwenden und weitere 16\,\% dies planen. 21,6\,\% haben den Einsatz
bereits diskutiert und nur 7,9\,\% geben an, dass \ac{ki} für sie keine Rolle spielt~\cite{e02}. Der deutlich höhere Anteil gegenüber den Zahlen des 
Statistischen Bundesamtes ist auf die unterschiedliche Erhebungsmethodik zurückzuführen, insbesondere auf eine abweichende Zusammensetzung der befragten Unternehmen
und eine weiter gefasste Definition der \ac{ki}-Nutzung. Beide Erhebungen zeigen jedoch übereinstimmend einen deutlichen Anstieg. \\
Mit dem zunehmenden Einsatz von \acp{llm} in Unternehmen greifen diese auch vermehrt auf Dokumente zu oder werden für weiterführende Aktionen wie Internetrecherchen 
eingesetzt. Dies spiegelt sich auch im Bericht der Organisation OWASP aus dem Jahr 2026 wider, in dem die zehn größten Risiken für \acp{llm} definiert werden. Darin wird
das Risiko \enquote{Excessive Agency} auf Platz~3 eingeordnet, während es 2025 noch auf Platz~6 lag~\cite{pi24}. Dieses Risiko beschreibt die Gefahr, die 
entsteht, wenn ein \ac{llm} auf zu viele Funktionalitäten zugreifen kann und über zu viele Berechtigungen verfügt. \\
Es zeigt sich somit, dass \acp{llm} zunehmend Zugriff auf große Datenmengen erhalten, auf deren Grundlage sie auch Aktionen ausführen können. Diese Aktionen werden als
Tool Calls bezeichnet. Durch Tool Calls wird es \acp{llm} ermöglicht, beliebige Funktionen auszuführen, beispielsweise auf externe Daten und Webseiten zuzugreifen 
oder Formulare auszufüllen. Es können unbegrenzt viele dieser Funktionen hinzugefügt werden. Mit dieser Erweiterung entstehen jedoch auch Gefahren, da durch den Aufruf 
externer Quellen und die Verarbeitung großer Datenmengen neue Angriffsflächen entstehen. Eine Ursache liegt darin, dass \acp{llm} nicht zuverlässig zwischen Anweisungen 
und Daten unterscheiden können. Dadurch können in den verarbeiteten Daten Anweisungen enthalten sein, die das \ac{llm} anschließend ausführt. Diese Schwachstelle
wird von Angreifern bei sogenannten Prompt-Injection-Angriffen ausgenutzt. Dabei werden Anweisungen in externen Daten platziert, auf die ein \ac{llm} zugreift.
Auf diese Weise lassen sich schädliche Aktionen auslösen, ohne dass ein direkter Zugriff auf das \ac{llm} notwendig ist~\cite{pi15}. \\
Durch den Einsatz von Tool Calls wird dieses Risiko verstärkt. Während ein erfolgreicher Angriff ohne Tool Call lediglich eine fehlerhafte Antwort erzeugt, können
mithilfe von Tool Calls reale Aktionen ausgeführt werden. Dadurch nimmt die Bedrohung dieser Angriffe deutlich zu. Besonders kritisch ist dies bei Chatbots, die 
auf Unternehmensdaten zugreifen, da diese häufig Zugang zu sensiblen internen Informationen haben. \\
Das Problem ist bereits seit mehreren Jahren bekannt, weshalb bereits zahlreiche Ansätze existieren, diese Angriffe abzuwehren. Es werden unter anderem zusätzliche
Klassifikatoren, Filter oder Vorverarbeitungsschritte implementiert, um Angriffe zu erkennen, bevor sie das \ac{llm} mit Zugriff auf die Tools erreichen~\cite[4]{pi9}~\cite[9]{sm10}.
Es gibt jedoch Belege dafür, dass diese Maßnahmen nicht ausreichend wirksam sind~\cite{sm22}. Angreifer entwickeln ihre Angriffe stetig weiter, um genau diese 
Schutzmechanismen zu umgehen. \\
Ein bekanntes Beispiel aus dem Jahr 2025 ist die Schwachstelle EchoLeak, welche Microsoft 365 Copilot betraf. Dabei gelang es Angreifern, mithilfe einer in einer E-Mail
versteckten Anweisung sensible Unternehmensdaten zu exfiltrieren und an einen eigenen Server zu übermitteln~\cite{e07}. Für diesen Angriff war keine Interaktion des
Nutzers erforderlich, sondern der Angriff wurde nur durch die E-Mail ausgelöst. EchoLeak stellt damit ein praxisnahes Beispiel für eine indirekte Prompt Injection
dar~\cite{e06}. \\
Einige der in dieser Arbeit berücksichtigten Arbeiten behandeln nicht ausschließlich indirekte Prompt-Injection-Angriffe, sondern teilweise Jailbreaks. Der Unterschied
besteht darin, dass sich Prompt Injection gegen Anwendungen richtet, die auf einem \ac{llm} aufbauen und bei denen nicht vertrauenswürdige Anweisungen mit einem 
vertrauenswürdigen Prompt verkettet werden. Bei einem Jailbreak wird hingegen versucht, die Sicherheitsfilter des Modells selbst zu umgehen. Das Merkmal, welches 
ausschlaggebend für Prompt Injection ist, ist somit die Verkettung eines vertrauenswürdigen mit einem nicht vertrauenswürdigen Inhalt kombiniert in einem Prompt~\cite{e05}. \\
Die Literatur zeigt, dass bei den bestehenden Schutzmaßnahmen noch erhebliches Verbesserungspotenzial besteht. An diesem Punkt setzt die vorliegende Arbeit an. Der 
Fokus liegt auf indirekten Prompt-Injection-Angriffen, bei denen der Angriff über externe Quellen erfolgt, auf die das \ac{llm} zugreift, beispielsweise eine Webseite~\cite[2\psq]{pi6}.
Ziel ist die Entwicklung einer Schutzmaßnahme, die genau diese Angriffe abwehrt und zugleich eine praxisnahe Anwendung ermöglicht und damit eine Anwendung bei verschiedenen
Chatbots. Die Schutzmaßnahme wird zunächst anhand eines selbst erstellten Datensatzes entwickelt, dessen Angriffe aus indirekten Angriffsszenarien abgeleitet sind, zur
Vereinfachung jedoch direkt im Nutzerprompt übergeben werden. Anschließend wird sie mit den öffentlichen Datensätzen BIPIA und NotInject überprüft. In dem Datensatz
BIPIA sind die Angriffe in externen Inhalten wie E-Mails, Tabellen oder Code positioniert. \\
Als Anwendungsfall dient der Chatbot ChatBot4U der Energy4U GmbH, der zudem als Unternehmenslösung für die Kommunikation mit Kunden vorgesehen ist und damit ein 
praxisnahes Anwendungsbeispiel darstellt. Im Folgenden werden die Zielsetzung, die Forschungsfragen sowie der Aufbau der Arbeit erläutert.

\section{Zielsetzung}

Ziel dieser Arbeit ist es zu untersuchen, inwieweit ein Chatbot, konkret der unternehmensinterne Chatbot der Energy4U GmbH, durch eine selbst entwickelte 
Schutzmaßnahme vor indirekten Prompt-Injection-Angriffen abgesichert werden kann. Als Vorstufe wird eine Schwachstellenanalyse durchgeführt, um die Angriffsfläche 
zu identifizieren. Hierfür wird eine lokale Testumgebung aufgebaut. \\
Auf Grundlage der identifizierten Schwachstellen wird eine geeignete Schutzmaßnahme entwickelt und implementiert. Diese wird anhand mehrerer Datensätze getestet und 
ihre Leistungsfähigkeit mithilfe geeigneter Metriken festgehalten. \\
Die Ergebnisse sollen einen direkten Vergleich zwischen einem einfachen vorgeschalteten \ac{llm}-Klassifikator, der jede Eingabe lediglich mit einer einzelnen Anfrage 
bewertet, und einer weiterentwickelten Schutzmaßnahme ermöglichen. Die Schutzmaßnahme soll sowohl für den unternehmensinternen Chatbot der Energy4U GmbH einsetzbar als 
auch auf andere praxisrelevante Chatbots übertragbar sein.

\section{Forschungsfragen}

\begin{enumerate}
    \item[HF:] Inwieweit kann ein unternehmensinterner Chatbot mit implementiertem Tool Calling gegen indirekte Prompt-Injection-Angriffe abgesichert werden?
    \item[TF1:] Welche Schwachstellen weist der Chatbot in Bezug auf indirekte Prompt-Injection-Angriffe auf?
    \item[TF2:] Wie zuverlässig ist ein vorgeschaltetes \ac{llm} bei der Abwehr von Prompt-Injection-Angriffen?
    \item[TF3:] Verbessert sich die Erkennung von Prompt-Injection-Angriffen, wenn die Schutzmaßnahme um Segmentierung und Rekonstruktion erweitert wird?
    \item[TF4:] Inwieweit lassen sich die Ergebnisse auf öffentliche Datensätze übertragen?
\end{enumerate}

Für diese Arbeit wurde eine Hauptforschungsfrage definiert, die in mehrere Teilfragen untergliedert ist. Die Hauptfrage lautet, inwieweit ein unternehmensinterner 
Chatbot mit implementiertem Tool Calling gegen indirekte Prompt-Injection-Angriffe abgesichert werden kann. Der unternehmensinterne Chatbot ist dabei das Testsystem.
Die Beantwortung der Frage soll allgemein gemacht werden. Zur Beantwortung der Hauptfrage wurden mehrere Teilfragen formuliert. Zunächst soll geklärt werden, welche 
Schwachstellen der Chatbot in Bezug auf indirekte Prompt-Injection-Angriffe aufweist. Während der Implementierung der Schutzmaßnahme wird untersucht, wie zuverlässig
ein vorgeschaltetes \ac{llm} Prompt-Injection-Angriffe abwehrt. Auf diesen Ergebnissen aufbauend soll geprüft werden, ob sich die Erkennung verbessert, wenn die 
Schutzmaßnahme in Kombination mit dem vorgeschalteten Modell um Rekonstruktionsmaßnahmen und eine Segmentierung der Daten erweitert wird. Abschließend wird untersucht,
inwieweit sich die Ergebnisse und der Ansatz auf öffentliche Datensätze übertragen lassen und wie praxisnah der Ansatz somit ist.

\section{Aufbau der Arbeit}

Die Arbeit gliedert sich in sieben Kapitel. Im Anschluss an die Einleitung werden im Grundlagenkapitel die fachlichen Grundlagen erarbeitet, um dem Leser einen 
Überblick über die Thematik zu verschaffen. Dazu werden \acp{llm} und ihre Funktionsweise, ihre Entwicklungsgeschichte, ihre technischen Bestandteile und 
ihr Aufbau behandelt. Darauf aufbauend wird Tool Calling bei \acp{llm} erläutert, also die Nutzung und der Ablauf eines Aufrufs von externen Ressourcen. Das Verständnis
dieser Funktion bildet eine wesentliche Grundlage für die Arbeit. Zusätzlich werden Sicherheitsbedrohungen vorgestellt, welche durch den Einsatz von Tool Calling 
entstehen. Anschließend werden Chatbots definiert und von \acp{llm} abgegrenzt. Es wird ihre Architektur erläutert und anhand dieser die Angriffsfläche beschrieben.
Abschließend werden unternehmensinterne Chatbots mit ihren Einsatzbereichen und Risiken betrachtet. \\
Im darauffolgenden Kapitel wird das zentrale Thema der Arbeit behandelt. Prompt-Injection-Angriffe werden definiert und in direkte und indirekte Angriffe kategorisiert.
Für beide Kategorien werden aktuelle Angriffe vorgestellt und erläutert. Zusätzlich wird die Bedeutung dieser Angriffe für Unternehmen behandelt und anhand eines
praxisnahen Beispiels veranschaulicht. \\
Auf Grundlage dieser Kapitel wird der praktische Teil der Arbeit vorbereitet, in dem der unternehmensinterne Chatbot auf Angriffe getestet wird. Hierfür wird eine 
Testumgebung aufgebaut. Es werden der Aufbau des Chatbots und seine Komponenten sowie die weiteren Bestandteile der Testumgebung beschrieben. Anschließend werden 
die Durchführung der einzelnen Angriffe beschrieben und die Ergebnisse dokumentiert, welcher Angriff in welchem Umfang erfolgreich war. Abschließend werden die aus
den Ergebnissen ableitbaren Schwachstellen dargestellt. \\
Im Anschluss an diese Vorarbeiten folgt der Hauptbestandteil der Arbeit: die Konzeption und Implementierung einer eigenen Schutzmaßnahme, mit der die zuvor dokumentierten
Angriffe abgewehrt werden sollen. Dazu werden zunächst mögliche Schutzmaßnahmen aus der Theorie aufbereitet und auf Basis ausgewählter Maßnahmen ein Konzept
entwickelt. Für die Schutzmaßnahme wird ein Datensatz erstellt sowie öffentliche Datensätze übertragen und geeignete Metriken ausgewählt. Anschließend wird die
Maßnahme implementiert und die Ergebnisse für die Auswertung aufbereitet. \\
In der Evaluation wird zunächst überprüft, ob die Forschungsfragen beantwortet wurden. Anschließend werden die Ergebnisse bewertet und analysiert, ob die Schutzmaßnahme
erfolgreich war und aus welchen Gründen. Zudem werden die Grenzen der Arbeit aufgezeigt und Verbesserungsmöglichkeiten definiert. \\
Abschließend werden die wesentlichen Ergebnisse der Arbeit zusammengetragen und ein Ausblick gegeben, wie zukünftige Forschung bestehende Lücken schließen kann.